% !TEX root = ../dabirestan.tex
%% ============================================================
%%  فصل ۲ — شمارش: اصل ضرب، جایگشت و ترکیب
%% ============================================================
\chapter{شمارش: اصل ضرب، جایگشت و ترکیب}
\label{hs:ch2}

\chapterintro{\begin{itemize}
  \item با اصل جمع و اصل ضرب، تعداد حالت‌ها را بدون فهرست کردن همه‌شان بشمارید؛
  \item جایگشت (چیدن با ترتیب) و ترکیب (انتخاب بدون ترتیب) را از هم تشخیص دهید؛
  \item با فرمول‌های $n!$، $P(n,r)$ و $\binom nr$ مسئله حل کنید.
\end{itemize}}

برای حساب کردن احتمال در فصل بعد باید بتوانیم تعداد برآمدها را بشماریم. وقتی تعداد حالت‌ها کم است فهرست کردنشان آسان است، ولی مثلاً تعداد رمزهای ۴ رقمی را نمی‌توانیم یکی‌یکی بنویسیم! در این فصل ابزارهای شمردن سریع را یاد می‌گیریم.

%% ------------------------------------------------------------
\section{اصل جمع و اصل ضرب}
%% ------------------------------------------------------------

\begin{theorem}[label=hs:thm:counting]{اصل جمع و اصل ضرب}
\begin{itemize}
  \item \textbf{اصل جمع:} اگر کاری را بتوان \emph{یا} به $m$ روش \emph{یا} به $n$ روش (بدون حالت مشترک) انجام داد، روی هم $m+n$ روش وجود دارد.
  \item \textbf{اصل ضرب:} اگر کاری دو مرحله داشته باشد و مرحلهٔ اول $m$ روش و مرحلهٔ دوم (برای هر انتخاب مرحلهٔ اول) $n$ روش داشته باشد، کل کار به $m\times n$ روش انجام می‌شود. برای چند مرحله هم تعدادها را در هم ضرب می‌کنیم.
\end{itemize}
\end{theorem}

\begin{example}[label=hs:ex:clothes]{لباس پوشیدن}
سارا ۳ پیراهن و ۲ شلوار دارد. به چند طریق می‌تواند یک دست لباس (یک پیراهن و یک شلوار) بپوشد؟
\begin{solution}
مرحلهٔ اول (پیراهن) ۳ انتخاب و مرحلهٔ دوم (شلوار) ۲ انتخاب دارد؛ طبق اصل ضرب $3\times2=6$ روش. نمودار درختی همین را نشان می‌دهد:
\begin{figblock}
\begin{tikzpicture}[
  grow=left, level distance=2.8cm,
  level 1/.style={sibling distance=1.9cm},
  level 2/.style={sibling distance=0.8cm},
  every node/.style={font=\small},
  st/.style={draw=primary!60, fill=primarylight, rounded corners=2pt, inner sep=3pt},
  edge from parent/.style={draw=primary!55, thick}]
  \node[circle, fill=primary, inner sep=2pt] {}
    child { node[st] {\rl{پیراهن ۱}} child { node {\rl{شلوار ۱}} } child { node {\rl{شلوار ۲}} } }
    child { node[st] {\rl{پیراهن ۲}} child { node {\rl{شلوار ۱}} } child { node {\rl{شلوار ۲}} } }
    child { node[st] {\rl{پیراهن ۳}} child { node {\rl{شلوار ۱}} } child { node {\rl{شلوار ۲}} } };
\end{tikzpicture}

\figcaption[hs:fig:clothes]{هر مسیر یک دست لباس است: $3\times2=6$ مسیر.}
\end{figblock}
\end{solution}
\end{example}

\begin{example}[label=hs:ex:passcode]{رمز گوشی}
یک رمز ۴ رقمی از ارقام ۰ تا ۹ ساخته می‌شود. (الف) چند رمز ممکن است؟ (ب) اگر رقم تکراری مجاز نباشد چطور؟
\begin{solution}
(الف) هر رقم ۱۰ انتخاب دارد: $10\times10\times10\times10=10\,000$.\\
(ب) رقم اول ۱۰ انتخاب، دومی ۹ (چون یکی استفاده شده)، سومی ۸ و چهارمی ۷: \ $10\times9\times8\times7=5040$.
\end{solution}
\end{example}

\begin{example}[label=hs:ex:menu]{اصل جمع}
در یک رستوران ۴ نوع غذای گوشتی و ۳ نوع غذای گیاهی هست. اگر بخواهید \emph{یک} غذا انتخاب کنید، چند انتخاب دارید؟ اگر بخواهید یک غذای گوشتی \emph{و} یک غذای گیاهی بگیرید چطور؟
\begin{solution}
یک غذا (گوشتی \emph{یا} گیاهی): طبق اصل جمع $4+3=7$. \ یکی از هر نوع (گوشتی \emph{و} گیاهی): طبق اصل ضرب $4\times3=12$.
\end{solution}
\end{example}

\begin{note}
یک قاعدهٔ سرانگشتی: «\textbf{یا}» معمولاً یعنی \textbf{جمع} و «\textbf{و} سپس» معمولاً یعنی \textbf{ضرب}.
\end{note}

%% ------------------------------------------------------------
\section{جایگشت: چیدن با ترتیب}
%% ------------------------------------------------------------

\begin{definition}[label=hs:def:factorial]{فاکتوریل}
برای عدد طبیعی $n$، \term{فاکتوریل}{factorial} $n$ عبارت است از
\[ n!=n\times(n-1)\times\cdots\times2\times1, \]
و قرارداد می‌کنیم $0!=1$. مثلاً $3!=6$، \ $4!=24$ و $5!=120$.
\end{definition}

\begin{theorem}[label=hs:thm:perm]{تعداد جایگشت‌ها}
\begin{itemize}
  \item $n$ شیء متمایز را به $n!$ روش می‌توان در یک ردیف چید.
  \item تعداد روش‌های انتخاب و چیدن $r$ شیء از بین $n$ شیء (با ترتیب) برابر است با
  \[ P(n,r)=n\times(n-1)\times\cdots\times(n-r+1)=\frac{n!}{(n-r)!}. \]
\end{itemize}
\end{theorem}

\textbf{چرا؟} طبق اصل ضرب: جایگاه اول $n$ انتخاب دارد، جایگاه دوم $n-1$ انتخاب، و همین‌طور تا آخر.

\begin{example}[label=hs:ex:books]{چیدن کتاب‌ها}
به چند طریق می‌توان ۵ کتاب متفاوت را در یک قفسه کنار هم چید؟
\begin{solution}
$5!=5\times4\times3\times2\times1=120$ روش.
\end{solution}
\end{example}

\begin{example}[label=hs:ex:podium]{سکوی مسابقه}
۸ دونده در یک مسابقه شرکت کرده‌اند. به چند طریق ممکن است مقام‌های اول، دوم و سوم مشخص شوند؟
\begin{solution}
ترتیب مهم است (اول بودن با دوم بودن فرق دارد): \ $P(8,3)=8\times7\times6=336$.
\end{solution}
\end{example}

%% ------------------------------------------------------------
\section{ترکیب: انتخاب بدون ترتیب}
%% ------------------------------------------------------------

گاهی فقط مهم است \emph{چه کسانی} انتخاب شوند، نه به چه ترتیبی. مثلاً انتخاب ۲ نفر از ۴ نفر برای یک تیم: تیم «علی و مریم» همان تیم «مریم و علی» است.

\begin{definition}[label=hs:def:comb]{ترکیب}
تعداد روش‌های انتخاب $r$ شیء از بین $n$ شیء متمایز \emph{بدون در نظر گرفتن ترتیب} را با $\binom nr$ (بخوانید «$n$ انتخاب $r$») نشان می‌دهیم و آن را \term{ترکیب}{combination} می‌نامیم:
\[ \binom nr=\frac{n!}{r!\,(n-r)!}. \]
\end{definition}

\textbf{چرا؟} اگر ترتیب مهم بود، $P(n,r)$ حالت داشتیم. ولی هر گروه $r$ نفره به $r!$ ترتیب مختلف در این شمارش آمده است؛ پس بر $r!$ تقسیم می‌کنیم: \ $\binom nr=\dfrac{P(n,r)}{r!}$.

\begin{example}[label=hs:ex:team]{انتخاب تیم}
از ۴ نفر (علی، مریم، رضا، سارا) می‌خواهیم ۲ نفر را برای یک تیم انتخاب کنیم. به چند طریق؟
\begin{solution}
$\binom42=\dfrac{4!}{2!\,2!}=\dfrac{24}{4}=6$. فهرستشان: علی‌مریم، علی‌رضا، علی‌سارا، مریم‌رضا، مریم‌سارا، رضاسارا. (اگر ترتیب مهم بود، $P(4,2)=12$ حالت داشتیم؛ هر تیم دو بار شمرده می‌شد.)
\end{solution}
\end{example}

\begin{example}[label=hs:ex:handshake]{دست دادن}
در یک جلسه ۱۰ نفر حاضرند و هر دو نفر یک بار با هم دست می‌دهند. چند دست دادن انجام می‌شود؟
\begin{solution}
هر دست دادن یعنی انتخاب ۲ نفر از ۱۰ نفر (بدون ترتیب): \ $\binom{10}2=\dfrac{10\times9}{2}=45$.
\end{solution}
\end{example}

\begin{example}[label=hs:ex:committee]{کمیته با شرط}
از ۵ دختر و ۴ پسر می‌خواهیم کمیته‌ای ۳ نفره بسازیم. (الف) به چند طریق؟ (ب) اگر کمیته دقیقاً ۲ دختر داشته باشد؟ (ج) اگر \emph{دست‌کم} یک پسر داشته باشد؟
\begin{solution}
(الف) $\binom93=84$.\\
(ب) ۲ دختر از ۵ و ۱ پسر از ۴: \ $\binom52\times\binom41=10\times4=40$.\\
(ج) شمردن «دست‌کم یک پسر» مستقیماً سخت است؛ از \emph{متمم} استفاده می‌کنیم: کمیته‌هایی که \emph{هیچ} پسری ندارند (هر سه دختر) $\binom53=10$ تا هستند؛ پس جواب $84-10=74$ است.
\end{solution}
\end{example}

\begin{center}
\renewcommand{\arraystretch}{1.5}
\begin{tabular}{l|c|c}
 & \textbf{ترتیب مهم است} & \textbf{ترتیب مهم نیست}\\\hline
\textbf{نام} & جایگشت & ترکیب\\
\textbf{فرمول} & $P(n,r)=\dfrac{n!}{(n-r)!}$ & $\dbinom nr=\dfrac{n!}{r!\,(n-r)!}$\\
\textbf{مثال} & رئیس، معاون و منشی & یک کمیتهٔ ۳ نفره\\
\end{tabular}
\end{center}

\begin{warning}
قبل از هر محاسبه از خودتان بپرسید: «\emph{اگر جای دو تا از انتخاب‌ها را عوض کنم، حالت جدیدی ساخته‌ام؟}» اگر بله، ترتیب مهم است (جایگشت)؛ اگر نه، ترتیب مهم نیست (ترکیب).
\end{warning}

%% ------------------------------------------------------------
\section*{خلاصهٔ فصل ۲}
\addcontentsline{toc}{section}{خلاصه و تمرین‌ها}
%% ------------------------------------------------------------

\begin{summary}[title={\leavevmode\boxicon{accent}{$\Sigma$}\ \ خلاصهٔ فصل}]
\begin{itemize}
  \item \textbf{اصل ضرب:} مراحل پشت سر هم $\Rightarrow$ ضرب. \ \textbf{اصل جمع:} حالت‌های جدا از هم $\Rightarrow$ جمع.
  \item \textbf{جایگشت} (ترتیب مهم): $n!$ و $P(n,r)=\frac{n!}{(n-r)!}$.
  \item \textbf{ترکیب} (ترتیب مهم نیست): $\binom nr=\frac{n!}{r!(n-r)!}$.
  \item برای «دست‌کم یکی» معمولاً راحت‌تر است حالت‌های «هیچ» را از کل کم کنیم.
\end{itemize}
\end{summary}

\subsection*{تمرین‌ها}

\begin{exercise}
پلاکی از ۲ حرف (از ۳۲ حرف الفبا) و سپس ۳ رقم تشکیل شده است. چند پلاک مختلف ممکن است؟
\answer{$32\times32\times1000=1\,024\,000$}
\end{exercise}

\begin{exercise}
به چند طریق می‌توان ۶ نفر را در یک صف ایستاد؟ اگر علی بخواهد حتماً اول صف باشد چطور؟
\answer{$720$؛ \ $120$}
\end{exercise}

\begin{exercise}
از یک کلاس ۲۵ نفره، یک نماینده و یک معاون انتخاب می‌شود. به چند طریق؟
\answer{$25\times24=600$}
\end{exercise}

\begin{exercise}
در یک آزمون باید به ۶ سؤال از ۸ سؤال جواب دهید. به چند طریق می‌توانید سؤال‌ها را انتخاب کنید؟
\answer{$\binom86=28$}
\end{exercise}

\begin{exercise}
یک جعبه ۶ مهرهٔ قرمز و ۴ مهرهٔ آبی دارد. به چند طریق می‌توان ۳ مهره برداشت که (الف) همه قرمز باشند؛ (ب) ۲ قرمز و ۱ آبی باشند؛ (ج) دست‌کم یک آبی باشد؟
\answer{(الف) $20$ \quad (ب) $60$ \quad (ج) $120-20=100$}
\end{exercise}

\begin{exercise}
با حروف کلمهٔ «کتاب» (۴ حرف متفاوت) چند کلمهٔ ۴ حرفی (با معنی یا بی‌معنی) می‌توان ساخت؟ چند کلمهٔ ۲ حرفی؟
\answer{$24$؛ \ $12$}
\end{exercise}
